% !TeX root = ../../Book4.tex
% 纲要标题：### A.2 单子上的代数与自由对象
% 纲要位置：工作记录/计划纲要.md，第 3828 行。
% 纲要细目（写作范围）：
% #### A.2.1 单子上的代数与同态
% #### A.2.2 自由代数与遗忘伴随
% #### A.2.3 比较函子与典型代数范畴


\section{单子上的代数与自由对象}
\label{b4:sec:A02}

单子给出形式运算，但形式表达式还须在一个对象中求值。求值要把单个元素还原为自身，并要求“先在内层求值再在外层求值”与“一次归并后求值”相同。单子上的代数就是满足这两个条件的对象。

\subsection{单子上的代数与同态}
\label{b4:subsec:A02:01}

\begin{definition}\label{b4:def:A-monad-algebra}
设 \((T,\eta,\mu)\) 为局部小范畴 \(\mathcal C\) 上的单子，\(X\in\mathcal C\)，\(a:TX\to X\) 为态射。若满足
\begin{equation}\label{b4:eq:A-algebra-laws}
a\eta_X=1_X,\qquad aT(a)=a\mu_X,
\end{equation}
则称 \((X,a)\) 为 \(T\) 上的\term[danzi-shang-daishu]{代数}[algebra for a monad]。对两个 \(T\)-代数 \((X,a)\)、\((Y,b)\)，若 \(\mathcal C\) 中的态射 \(f:X\to Y\) 满足
\[
fa=bT(f)
\]
则称 \(f\) 为它们之间的代数同态。所有这些代数与同态组成的范畴记为 \(\mathcal C^T\)，称为 \term[eilenberg-moore-fanchou]{Eilenberg–Moore 范畴}[Eilenberg--Moore category]。\footnote{摩尔（John Coleman Moore，1923-2016），美国数学家，研究代数拓扑和谱序列；艾伦伯格已见本卷前注。}
\end{definition}
\symidx{\(\mathcal C^T\)：单子 \(T\) 上的代数范畴}

恒等态射满足同态等式；若 \(f a=bT(f)\)、\(g b=cT(g)\)，则
\((gf)a=cT(gf)\)。所以它们确实组成范畴。遗忘函子
\(U^T:\mathcal C^T\to\mathcal C\) 把 \((X,a)\) 送到 \(X\)，把同态送到底层态射；它是忠实函子，但一般不是全函子。

\begin{proposition}\label{b4:prop:A-idempotent-algebras}
设 \((T,\eta,\mu)\) 为局部小范畴 \(\mathcal C\) 上的幂等单子，\(X\in\mathcal C\)。\(X\) 存在 \(T\)-代数结构，当且仅当单位分量 \(\eta_X:X\to TX\) 可逆；此时结构唯一，为 \(a=\eta_X^{-1}\)。满足此条件的对象之间的每个 \(\mathcal C\)-态射都是代数同态。
\end{proposition}
\begin{proof}
若 \(a\eta_X=1_X\)，由 \(\eta\) 的自然性及幂等性得到
\[
\eta_Xa=T(a)\eta_{TX}=T(a)T(\eta_X)=1_{TX}.
\]
所以 \(a\) 必为 \(\eta_X\) 的逆。反之，若 \(\eta_X\) 可逆，令 \(a=\eta_X^{-1}\)。单位律显然成立，且
\[
T(a)=T(\eta_X)^{-1}=\mu_X,
\]
从而 \(aT(a)=a\mu_X\)。若 \(f:X\to Y\)，自然性 \(T(f)\eta_X=\eta_Yf\) 在两边乘以单位的逆，就得到代数同态等式。
\end{proof}

\subsection{自由代数与遗忘伴随}
\label{b4:subsec:A02:02}

\begin{theorem}[自由代数与遗忘伴随]\label{b4:thm:A-em-adjunction}
设 \((T,\eta,\mu)\) 为局部小范畴 \(\mathcal C\) 上的单子，\(\mathcal C^T\) 为其代数范畴，\(U^T:\mathcal C^T\to\mathcal C\) 为取底层对象与态射的遗忘函子。对 \(X\in\mathcal C\) 及 \(\mathcal C\) 中的态射 \(f:X\to X'\)，公式
\[
F^T(X)=(TX,\mu_X),\qquad F^T(f)=T(f)
\]
定义函子 \(F^T:\mathcal C\to\mathcal C^T\)，且 \(F^T\dashv U^T\)。对任意 \(T\)-代数 \((Y,b)\)，伴随双射为
\begin{equation}\label{b4:eq:A-em-adjunction}
\operatorname{Hom}_{\mathcal C^T}\bigl(F^TX,(Y,b)\bigr)
\cong\operatorname{Hom}_{\mathcal C}(X,Y),
\qquad h\longmapsto h\eta_X,
\end{equation}
其逆把 \(f\) 送到 \(bT(f)\)。这个伴随产生的单子正是原来的 \(T\)。
\end{theorem}
\begin{proof}
\eqref{b4:eq:A-monad-laws}说明 \((TX,\mu_X)\) 满足代数公理。乘法的自然性
\(T(f)\mu_X=\mu_YT^2(f)\) 说明 \(T(f)\) 是代数同态，函子性来自 \(T\)。

对 \(f:X\to Y\)，令 \(\overline f=bT(f)\)。利用代数公理及乘法自然性，
\[
\overline f\mu_X
=b\mu_YT^2(f)
=bT(b)T^2(f)=bT(\overline f),
\]
所以它是代数同态。又
\(\overline f\eta_X=b\eta_Yf=f\)。
反过来，若 \(h:(TX,\mu_X)\to(Y,b)\) 为代数同态，则
\[
bT(h\eta_X)=bT(h)T(\eta_X)
=h\mu_XT(\eta_X)=h.
\]
故两公式互逆。与源、目标的态射复合时，这两个公式仍相容，因而构成自然双射。

单位为 \(\eta_X\)，余单位在 \((Y,b)\) 处就是 \(b:(TY,\mu_Y)\to(Y,b)\)。所以遗忘后复合自由函子仍为 \(T\)，伴随产生的乘法在 \(X\) 处为自由代数的结构态射 \(\mu_X\)，与原数据相同。
\end{proof}

例如，有限表单子的代数给出幺半群。对求值 \(a:TX\to X\)，定义
\[
e=a\bigl([]\bigr),\qquad x\cdot y=a\bigl([x,y]\bigr).
\]
对由三张单元素表组成的分组，以及把其中两张先连接的分组，代数结合律说明
\((x\cdot y)\cdot z=x\cdot(y\cdot z)=a\bigl([x,y,z]\bigr)\)；对含空表的分组得到左右单位律。归纳可知 \(a\bigl([x_1,\ldots,x_n]\bigr)=x_1\cdots x_n\)。反过来，幺半群的有序乘积满足两条代数公理。代数同态恰好保持空表和二元素表的值，也就是幺半群同态。

\ProseParagraphBreak

对任意 \(T\)-代数 \((Y,b)\)，单位律 \(b\eta_Y=1_Y\) 说明余单位 \(b\) 在底层范畴 \(\mathcal C\) 中是分裂满态射，但这并不自动给出代数范畴中的分裂。有限表单子把幺半群 \(Y\) 的表送到乘积，底层截面为 \(y\mapsto[y]\)；然而
\[
\eta_Y(y_1y_2)=[y_1y_2],\qquad
\eta_Y(y_1)\eta_Y(y_2)=[y_1,y_2],
\]
两张表不同，所以这个截面不是幺半群同态。某些求值映射甚至不存在任何幺半群同态截面：取二元群 \(Y=\{e,g\}\)，其中 \(g^2=e\)。若有幺半群同态 \(\sigma:Y\rightarrow TY\)，则 \(\sigma(g)\sigma(g)=\sigma(e)=[]\)，从表长相加可知 \(\sigma(g)=[]\)。于是 \(b\sigma(g)=e\ne g\)，故 \(\sigma\) 不可能是 \(b\) 的截面。要从自由代数恢复一般代数，还需把求值所满足的关系纳入构造；下一节将用余等化子完成这一步。

\subsection{比较函子与典型代数范畴}
\label{b4:subsec:A02:03}

\begin{proposition}\label{b4:prop:A-comparison-functor}
设 \(F:\mathcal C\to\mathcal D\)、\(U:\mathcal D\to\mathcal C\) 为局部小范畴之间的伴随 \(F\dashv U\)，单位、余单位为 \(\eta,\varepsilon\)，所生单子为 \(T=UF\)。则
\[
K(Y)=\bigl(UY,U(\varepsilon_Y)\bigr),\qquad K(f)=U(f)
\]
定义函子 \(K:\mathcal D\to\mathcal C^T\)，满足 \(U^TK=U\)、\(KF=F^T\)。它称为该伴随的\term[bijiao-hanzi]{比较函子}[comparison functor]。
\end{proposition}
\begin{proof}
记 \(a_Y=U(\varepsilon_Y)\)。三角恒等式给出 \(a_Y\eta_{UY}=1_{UY}\)。把余单位自然性用于 \(\varepsilon_Y:FUY\to Y\)，得到
\[
\varepsilon_YFU(\varepsilon_Y)
=\varepsilon_Y\varepsilon_{FUY}.
\]
应用 \(U\)，便有 \(a_YT(a_Y)=a_Y\mu_{UY}\)，所以 \(K(Y)\) 是代数。对 \(f:Y\to Z\)，余单位自然性又给出
\(U(f)a_Y=a_ZT\bigl(U(f)\bigr)\)，因此 \(K(f)\) 是代数同态。其余函子等式来自定义。
\end{proof}

等式 \(KF=F^T\) 和 \(U^TK=U\) 表明，原伴随的自由对象与单子重新构造的自由代数总能对应，且比较函子保持底层对象与态射。问题是原范畴 \(\mathcal D\) 中的全部对象和态射能否由这些代数数据恢复。为此需要检查 \(K\) 是否为范畴等价：每个 \(T\)-代数都应同构于某个 \(K(Y)\)，而 \(Y,Z\) 之间的态射也应与 \(K(Y),K(Z)\) 之间的代数同态一一对应。仅有自由伴随 \(F^T\dashv U^T\) 的存在还不能保证这两件事。

\begin{definition}\label{b4:def:A-monadic-functor}
设 \(\mathcal C,\mathcal D\) 为局部小范畴，函子 \(U:\mathcal D\to\mathcal C\) 有左伴随 \(F:\mathcal C\to\mathcal D\)，伴随的余单位为 \(\varepsilon:FU\Rightarrow1_{\mathcal D}\)。以 \(\mathcal C^{UF}\) 表示所生单子的代数范畴，定义比较函子 \(K(Y)=(UY,U(\varepsilon_Y))\)、\(K(f)=U(f)\)。若 \(K:\mathcal D\to\mathcal C^{UF}\) 是范畴等价，则称 \(U\) 为\term[danzixing-hanzi]{单子性函子}[monadic functor]，也说这个伴随具有单子性。
\end{definition}

对自由幺半群与底集合的伴随，比较函子把幺半群送到“有限表取乘积”的代数。上一小节已经证明每个代数都唯一确定这种乘法，而且代数同态就是幺半群同态，所以遗忘函子具有单子性。\(\mathcal C^T\) 在这里保留了全部代数结构。

\ProseParagraphBreak

并非每个伴随都能如此恢复另一范畴。离散拓扑函子
\(F:\mathbf{Set}\to\mathbf{Top}\) 左伴随于底集合函子 \(U\)，且 \(UF=1_{\mathbf{Set}}\)，所生单子只是恒等单子。它的代数范畴仍为集合范畴，不能恢复所有拓扑。具体地，同一个二点集从离散拓扑到平凡拓扑的恒等映射是连续双射，却不是同胚；下节的判据会指出这个例子在哪项假设上失败。
\symidx{\(\mathbf{Top}\)：拓扑空间及连续映射组成的范畴}
